\scanpage{44}\mspage{43}
\uline{Cont 8} Soit $X,Y,X'$ dans $M$, $Y\ll X$, $X'<X$. Considérons l'ens des $Y'\in\deleted{ M}F_Y$ (i.e. $Y'\leq Y$) tels que $Y'\ll X$ :
% slanted
\margin{Axiome du raffinement induit.}
\diagram{\begin{tikzcd}[ampersand replacement=\&,row sep=0.1em,column sep=0.35em]
Y \& \ll \& X\\
\tikz[baseline=-.5ex]\node[rotate=90,inner sep=0pt] {$\leq$}; \& {} \& \tikz[baseline=-.5ex]\node[rotate=90,inner sep=0pt] {$\unsure{<|\leq}$};\\
Y' \& \ll \& X'
\end{tikzcd}}

(lesquels forment une sous-figure $R_{X'}$ de le multistrate $\deleted{\unsure{F|\in}}F_Y$, \deleted{qu'on appellera} raffinant $X'$, et qu'on appellera le raffinement de $X'$, induit par $Y$ sur $X'$). Alors \unsure{tout\deleted{e}|toute} raffinement commun $Z$ de $Y$ et de $X'$ raffine un des $Y'$, i.e. raffine $R_{X'}$. (Donc Tout figure qui raffine à la fois $\deleted{\unsure{Y|F}}F_Y$, et $F_{X'}$, raffine $R_{X'}$, qui est \deleted{\unsure{donc} la borne} borne inf. de $F_Y$ et $F_{X'}$ pour la relation $\ll$) En particulier,
\[
 :\quad \widetilde Y\cap\widetilde{X'}=\bigcup_{Y'\in R_{X'}}\widetilde{Y'}
\]

\uline{Cor} \deleted{\unsure{Soit} $x\in\widetilde{\unsure{Y|X}}^\circ$,} Supposons que $\deleted{\unsure{\exists}}\nexists\,X_1\in M$, avec $Y\ll X_1<X$ (c'est la situation de tantôt) Alors $\widetilde Y^\circ\subset\widetilde X^\circ$.

En effet, si $x\in\widetilde Y^\circ$, \unsure{si on|s'on} avait $x\in\widetilde{X'}$ pour $X'<X$, on aurait $x\in\widetilde{Y'}$ avec $Y'\in R_{X'}$, donc $Y'\Rightarrow\deleted{\unsure{X}}Y$ par déf. de $\widetilde Y^\circ$, donc $Y\ll X'_{\deleted{\unsure{1}}}<X$, \inserted{contraire à hyp.}
